\clearpage
\chapter*{GLOSARIO}
\addcontentsline{toc}{chapter}{GLOSARIO}

\noindent\textit{Definición de los términos técnicos y acrónimos relevantes utilizados a lo largo del documento:}

\begin{description}
  \item[API (Application Programming Interface)] Interfaz que permite la comunicación entre diferentes componentes de software, utilizada en este proyecto para conectar servicios externos.
  \item[Burbujas Flotantes] Elementos de interfaz de usuario superpuestos dinámicamente sobre el canvas WebGL para presentar sugerencias de interacción de forma amigable.
  \item[Chunking] Proceso de dividir o fraccionar información. En este documento refiere al particionamiento de texto para indexación (RAG) o la transmisión de audio fragmentado (\textit{Streaming}).
  \item[Cosine Similarity (Similitud Coseno)] Métrica matemática utilizada para medir la semejanza entre dos vectores en un espacio multidimensional, fundamental para la recuperación de información.
  \item[Droplet] Servidor virtual privado escalable provisto por la infraestructura en la nube de DigitalOcean.
  \item[ECA (Embodied Conversational Agent)] Agente conversacional que cuenta con una representación corporal visible, como un avatar 3D.
  \item[FastAPI] Marco de trabajo (\textit{framework}) moderno y de alto rendimiento para construir APIs con Python.
  \item[GPU (Graphics Processing Unit)] Unidad de procesamiento especializada, crucial para la aceleración de modelos de inteligencia artificial y el renderizado gráfico.
  \item[LLM (Large Language Model)] Modelo de Inteligencia Artificial entrenado con cantidades masivas de texto capaz de comprender y generar lenguaje natural.
  \item[OCR (Optical Character Recognition)] Tecnología utilizada para extraer texto desde imágenes o documentos escaneados.
  \item[PLN (Procesamiento de Lenguaje Natural)] Rama de la inteligencia artificial enfocada en la interacción entre las computadoras y el lenguaje humano.
  \item[RAG (Retrieval-Augmented Generation)] Arquitectura de IA que mejora las respuestas de un LLM al recuperar previamente información relevante de una base de datos de conocimiento específica.
  \item[Three.js] Biblioteca de JavaScript que facilita la creación y visualización de gráficos 3D en el navegador web mediante WebGL.
  \item[TTS (Text-to-Speech)] Tecnología capaz de convertir texto escrito a audio sintetizado.
  \item[VRM (Virtual Reality Model)] Formato de archivo estándar diseñado específicamente para la representación de avatares 3D interoperables en aplicaciones de realidad virtual o navegadores.
  \item[WebSockets] Protocolo de comunicación bidireccional, simultánea y en tiempo real a través de una única conexión TCP.
  \item[Whisper] Modelo de reconocimiento de voz y traducción (\textit{Speech-to-Text}) de código abierto desarrollado por OpenAI.
\end{description}